\section{Robustness and sensitivity analysis}
\label{sec:robustness}

The main specification combines a rich article representation with a
particular temporal aggregation rule. This section asks how forecast
accuracy changes when those components are varied. We focus on BTC
four-hour forecasts and use all six expanding windows. Each experiment
loads the same fitted first-phase TS-only model used in the main
results, freezes its weights, and estimates a new second-phase news
correction. The learning rate, loss, batch size, epoch limit, early
stopping rule, and gradient clipping match the main NED specification.
Thus, a difference between variants reflects the fitted news
correction rather than a different first-phase baseline.

We refit the full-dimensional exponential specification along with
the alternatives. This provides a common reference trained with the
same numerical implementation and shuffled batch sequence as the
ablation models. It need not reproduce the main-results checkpoint
exactly, since the latter uses a different batch-generation path.
The experiments use one fixed training seed and report arithmetic
means across the six test windows. Their differences are descriptive
comparisons rather than estimates averaged over training seeds.

\subsection{Alternative decay specifications}
\label{sec:decay_ablation}

We first replace the exponential kernel while retaining the
768-dimensional input, LayerNorm, article head, residual map, and
twelve-hour lookback. If $u_k$ is the second unconstrained output
of the article head, the four article signals are
\begin{align*}
 \text{Constant:}\quad &s_k,\\
 \text{Power-law:}\quad
 &s_k(1+\Delta t_k/60)^{-\softplus(u_k)},\\
 \text{Exponential:}\quad
 &s_k\exp\{-\lambda_k\Delta t_k\},
 \qquad \lambda_k=0.001+0.499\sigmoid(u_k),\\
 \text{Linear:}\quad
 &s_k\max\!\left\{0,1-\frac{\Delta t_k}{1+60\softplus(u_k)}\right\}.
\end{align*}
Article age $\Delta t_k$ is measured in minutes. The power-law
specification fixes its time scale at sixty minutes and learns a
positive exponent. The linear specification learns a positive
duration exceeding one minute. All signals are truncated at the
lookback boundary. The constant variant leaves its second output
unused and provides no within-window discounting.

\begin{table}[htbp]
\centering\small
\caption{Alternative decay specifications (4-Hour RV, BTC)}
\label{tab:decay_ablation}
\begin{adjustbox}{max width=\textwidth}
\begin{tabular}{lrrrr}
\toprule
Decay function & $R^2$ & MSE & QLIKE & DA \\
\midrule
TS-only (no news) & 0.419 & 0.810 & 0.638 & 0.717 \\
Constant (no decay) & 0.423\,$\pm$0.035 & 0.805\,(-0.6\%) & 0.592\,(-7.2\%) & 0.718 \\
Power-law & 0.433\,$\pm$0.038 & 0.791\,(-2.3\%) & 0.568\,(-10.9\%) & 0.722 \\
Exponential & 0.434\,$\pm$0.040 & 0.790\,(-2.5\%) & 0.571\,(-10.6\%) & 0.723 \\
Linear & 0.440\,$\pm$0.037 & 0.781\,(-3.6\%) & 0.549\,(-14.0\%) & 0.725 \\
\bottomrule
\end{tabular}

\end{adjustbox}
\tabnote{Results are averaged across six expanding windows.
$\pm$ denotes the sample standard deviation of $R^2$ across windows.
MSE and $R^2$ are calculated on log-RV; QLIKE and DA use levels.
Parentheses give percentage changes relative to TS-only, with
negative values indicating lower loss. Every news specification,
including exponential, is fitted in the matched ablation experiment.}
\end{table}

Table~\ref{tab:decay_ablation} shows that the constant kernel already
improves the shared TS-only baseline. Its mean QLIKE is 0.59249,
a reduction of 7.16\%, and its mean log-MSE is 0.80461,
a reduction of 0.63\%. A model without within-window discounting
can therefore improve forecast accuracy. The evidence does not
support treating temporal decay as a necessary condition for any
benefit from the news correction.

All three decaying specifications improve further on both average
losses. The refitted exponential model reduces QLIKE by 10.57\%
and log-MSE by 2.48\% relative to TS-only. Power-law decay yields
a slightly larger QLIKE reduction, 10.94\%, and a slightly smaller
log-MSE reduction, 2.31\%. Linear decay has the lowest losses in
this comparison: QLIKE is 0.54907 and log-MSE is 0.78077,
corresponding to reductions of 13.96\% and 3.58\%. It also has
the highest mean $R^2$, 0.440, and DA, 0.725.

The comparison supports a role for temporal discounting beyond a
constant aggregation of article magnitudes, while qualifying a
stronger claim that the gain disappears without decay. It also
shows that functional form can matter quantitatively. Exponential
decay offers a smooth response and a direct half-life interpretation,
but it is not the most accurate kernel in this experiment. Since
each specification is estimated with one fixed seed, these rankings
describe the fitted models rather than establishing an ordering
that is invariant to optimization variability.


\subsection{Alternative news representations}
\label{sec:news_representations}

The second exercise varies the information supplied to the news
correction. Panel A of Table~\ref{tab:news_representations} uses
deterministic per-bar news aggregates without learned article decay.
These are the log of one plus the twelve-hour article count,
sentiment probabilities averaged with weights proportional to
$1/(1+\Delta t_k)$, sentiment probabilities averaged over the current
UTC calendar day up to the forecast cutoff, and the mean
768-dimensional CryptoBERT embedding over the preceding twelve hours.
The daily aggregate uses only articles already published; it is
not an end-of-day sentiment measure.

For these variants, a small residual head maps the bar-level news
feature into a log-variance correction, clipped to $[-2,2]$.
The count uses a linear map; the other inputs use a hidden layer
of width sixteen. Except for the log-count, features are standardized
using training-period moments. Empty sets are represented by zero
before standardization. The price branch remains frozen, as in the
article-decay models.

Panel B retains the learned exponential article-decay mechanism
and changes its input to the three FinBERT class probabilities,
a 256-dimensional PCA projection of CryptoBERT, or the full
768-dimensional CryptoBERT embedding. PCA is fitted separately in
each expanding window using retained training-period articles only.
All three article heads have hidden widths 128 and 32; LayerNorm
is applied to the embedding inputs and omitted for sentiment
probabilities. Thus the sentiment comparison uses a three-dimensional
probability vector rather than a single signed score.

\begin{table}[htbp]
\centering\small
\caption{Alternative news representations (4-Hour RV, BTC)}
\label{tab:news_representations}
\begin{adjustbox}{max width=\textwidth}
\begin{tabular}{lrrr}
\toprule
Representation & $R^2$ & MSE & QLIKE \\
\midrule
\multicolumn{4}{@{}l}{\textit{A. Without learned decay}} \\[3pt]
TS-only (no news) & 0.419 & 0.810 & 0.638 \\
+ News count & 0.421 & 0.807\,(-0.4\%) & 0.602\,(-5.7\%) \\
+ Sentiment $\times$ recency & 0.420 & 0.808\,(-0.2\%) & 0.598\,(-6.3\%) \\
+ Daily agg.\ sentiment & 0.431 & 0.793\,(-2.0\%) & 0.581\,(-9.0\%) \\
+ CryptoBERT 768-dim mean-pool & 0.422 & 0.805\,(-0.5\%) & 0.596\,(-6.6\%) \\
\addlinespace\multicolumn{4}{@{}l}{\textit{B. With learned $(s,\lambda)$ decay}} \\[3pt]
+ Sentiment & 0.433 & 0.791\,(-2.3\%) & 0.570\,(-10.7\%) \\
+ CryptoBERT 256-dim (PCA) & 0.437 & 0.785\,(-3.0\%) & 0.558\,(-12.6\%) \\
+ CryptoBERT 768-dim & 0.434 & 0.790\,(-2.5\%) & 0.571\,(-10.6\%) \\
\bottomrule
\end{tabular}

\end{adjustbox}
\tabnote{Results are arithmetic means across six windows. Parentheses
show loss changes relative to the shared TS-only baseline (negative
indicates improvement). Panel A uses bar-level news residuals without
learned article decay. Panel B learns article-level $(s,\lambda)$
and an aggregate residual map. All variants use frozen first-phase
price weights. The full-dimensional entry is the same refitted
exponential model as in Table~\ref{tab:decay_ablation}.}
\end{table}

Panel A shows that simple news aggregates can already improve the
time-series forecast. News count, recency-weighted sentiment, and
mean-pooled CryptoBERT reduce QLIKE by 5.74\%, 6.32\%, and 6.64\%,
respectively, although their log-MSE gains are below 0.6\%.
The current-day sentiment average performs better on both losses,
reducing QLIKE by 9.03\% and log-MSE by 2.01\%. Thus a useful
news correction does not require a learned article-decay kernel,
and a higher-dimensional aggregate is not automatically more useful
than a compact sentiment summary.

Panel B shows that sentiment combined with learned decay nearly
matches the full-dimensional embedding. Its QLIKE reduction is
10.68\%, compared with 10.57\% for the refitted CryptoBERT model;
its log-MSE reduction is 2.34\%, compared with 2.48\%.
These are small numerical differences with opposite signs across
the two losses. The sentiment probability vector therefore captures
most of the improvement in the full-dimensional reference, without
establishing that the representations are equivalent in other
samples or training runs.

The PCA representation has the lowest average losses within this
table. With 256 components, QLIKE falls by 12.55\% and log-MSE
by 3.02\% relative to TS-only, exceeding the reductions from the
768-dimensional input. Compression consequently does not sacrifice
forecast accuracy in this experiment. The result also cautions
against interpreting embedding dimension as a direct measure of
the useful information delivered to the news head.

Together, the comparisons support learning how article signals
change with age, but do not place all of the forecast benefit in
that mechanism. Deterministic aggregation already helps, while
sentiment and compressed embeddings remain competitive when learned
decay is used. Panel A and Panel B also differ in aggregation and
head structure, so their contrast does not isolate encoder quality.
The evidence favors evaluating representation and temporal weighting
jointly, with the matched constant-kernel comparison providing the
more direct test of within-window discounting.

